% OpenStax Calculus Volume 2 -- Section 3.7 Exercises: Improper Integrals
% Exercises reproduced from OpenStax, licensed under CC BY 4.0.
% Source: https://openstax.org/books/calculus-volume-2/pages/3-7-improper-integrals
\documentclass{exercises}
\title{OpenStax Calculus Volume 2}
\subtitle{Section 3.7 Exercises: Improper Integrals}
\sourceurl{https://openstax.org/books/calculus-volume-2/pages/3-7-improper-integrals}
\author{}
\date{}

\begin{document}
\maketitle


Evaluate the following integrals. If the integral is not convergent, answer ``divergent.''

\begin{enumerate}[start=1]
\item $\int_{2}^{4}\frac{dx}{{(x-3)}^{2}}$
\item $\int_{0}^{\infty}\frac{1}{4+x^{2}}\,dx$
\item $\int_{0}^{2}\frac{1}{\sqrt{4-x^{2}}}\,dx$
\item $\int_{1}^{\infty}\frac{1}{x\ln x}\,dx$
\item $\int_{1}^{\infty}xe^{-x}\,dx$
\item $\int_{-\infty}^{\infty}\frac{x}{x^{2}+1}\,dx$
\item Without integrating, determine whether the integral $\int_{1}^{\infty}\frac{1}{\sqrt{x^{3}+1}}\,dx$ converges or diverges by comparing the function $f(x)=\frac{1}{\sqrt{x^{3}+1}}$ with $g(x)=\frac{1}{\sqrt{x^{3}}}$.
\item Without integrating, determine whether the integral $\int_{1}^{\infty}\frac{1}{\sqrt{x+1}}\,dx$ converges or diverges.
\end{enumerate}

Determine whether the improper integrals converge or diverge. If possible, determine the value of the integrals that converge.

\begin{enumerate}[resume]
\item $\int_{0}^{\infty}e^{-x}\cos x\,dx$
\item $\int_{1}^{\infty}\frac{\ln x}{x}\,dx$
\item $\int_{0}^{1}\frac{\ln x}{\sqrt{x}}\,dx$
\item $\int_{0}^{1}\ln x\,dx$
\item $\int_{-\infty}^{\infty}\frac{1}{x^{2}+1}\,dx$
\item $\int_{1}^{5}\frac{dx}{\sqrt{x-1}}$
\item $\int_{-2}^{2}\frac{dx}{{(1+x)}^{2}}$
\item $\int_{0}^{\infty}e^{-x}\,dx$
\item $\int_{0}^{\infty}\sin x\,dx$
\item $\int_{-\infty}^{\infty}\frac{e^{x}}{1+e^{2x}}\,dx$
\item $\int_{0}^{1}\frac{dx}{\sqrt[3]{x}}$
\item $\int_{0}^{2}\frac{dx}{x^{3}}$
\item $\int_{-1}^{2}\frac{dx}{x^{3}}$
\item $\int_{0}^{1}\frac{dx}{\sqrt{1-x^{2}}}$
\item $\int_{0}^{3}\frac{1}{x-1}\,dx$
\item $\int_{1}^{\infty}\frac{5}{x^{3}}\,dx$
\item $\int_{3}^{5}\frac{5}{{(x-4)}^{2}}\,dx$
\end{enumerate}

Determine the convergence of each of the following integrals by comparison with the given integral. If the integral converges, find the number to which it converges.

\begin{enumerate}[resume]
\item $\int_{1}^{\infty}\frac{dx}{x^{2}+4x}$; compare with $\int_{1}^{\infty}\frac{dx}{x^{2}}$.
\item $\int_{1}^{\infty}\frac{dx}{\sqrt{x}+1}$; compare with $\int_{1}^{\infty}\frac{dx}{2\sqrt{x}}$.
\end{enumerate}

Evaluate the integrals. If the integral diverges, answer ``diverges.''

\begin{enumerate}[resume]
\item $\int_{1}^{\infty}\frac{dx}{x^{e}}$
\item $\int_{0}^{1}\frac{dx}{x^{\pi}}$
\item $\int_{0}^{1}\frac{dx}{\sqrt{1-x}}$
\item $\int_{0}^{1}\frac{dx}{1-x}$
\item $\int_{-\infty}^{0}\frac{dx}{x^{2}+1}$
\item $\int_{-1}^{1}\frac{dx}{\sqrt{1-x^{2}}}$
\item $\int_{0}^{1}\frac{\ln x}{x}\,dx$
\item $\int_{0}^{e}\ln(x)\,dx$
\item $\int_{0}^{\infty}xe^{-x}\,dx$
\item $\int_{-\infty}^{\infty}\frac{x}{{(x^{2}+1)}^{2}}\,dx$
\item $\int_{0}^{\infty}e^{x}\,dx$
\end{enumerate}

Evaluate the improper integrals. Each of these integrals has an infinite discontinuity either at an endpoint or at an interior point of the interval.

\begin{enumerate}[resume]
\item $\int_{0}^{9}\frac{dx}{\sqrt{9-x}}$
\item $\int_{-27}^{1}\frac{dx}{x^{2/3}}$
\item $\int_{0}^{3}\frac{dx}{\sqrt{9-x^{2}}}$
\item $\int_{6}^{24}\frac{dt}{t\sqrt{t^{2}-36}}$
\item $\int_{0}^{4}x\ln(4x)\,dx$
\item $\int_{0}^{3}\frac{x}{\sqrt{9-x^{2}}}\,dx$
\item Evaluate $\int_{.5}^{1}\frac{dx}{\sqrt{1-x^{2}}}$. (Be careful!) (Express your answer using three decimal places.)
\item Evaluate $\int_{1}^{4}\frac{dx}{\sqrt{x^{2}-1}}$. (Express the answer in exact form.)
\item Evaluate $\int_{2}^{\infty}\frac{dx}{{(x^{2}-1)}^{3/2}}$.
\item Find the area of the region in the first quadrant between the curve $y=e^{-6x}$ and the $x$-axis.
\item Find the area of the region bounded by the curve $y=\frac{7}{x^{2}}$, the $x$-axis, and on the left by $x=1$.
\item Find the area under the curve $y=\frac{1}{{(x+1)}^{3/2}}$, bounded on the left by $x=3$.
\item Find the area under $y=\frac{5}{1+x^{2}}$ in the first quadrant.
\item Find the volume of the solid generated by revolving about the $x$-axis the region under the curve $y=\frac{3}{x}$ from $x=1$ to $x=\infty$.
\item Find the volume of the solid generated by revolving about the $y$-axis the region under the curve $y=6e^{-2x}$ in the first quadrant.
\item Find the volume of the solid generated by revolving about the $x$-axis the area under the curve $y=3e^{-x}$ in the first quadrant.
\end{enumerate}

The Laplace transform of a continuous function over the interval $[0,\infty)$ is defined by $F(s)=\int_{0}^{\infty}e^{-sx}f(x)\,dx$ (see the Student Project). This definition is used to solve some important initial-value problems in differential equations, as discussed later. The domain of $F$ is the set of all real numbers $s$ such that the improper integral converges. Find the Laplace transform $F$ of each of the following functions and give the domain of $F$.

\begin{enumerate}[resume]
\item $f(x)=1$
\item $f(x)=x$
\item $f(x)=\cos(2x)$
\item $f(x)=e^{ax}$
\item Use the formula for arc length to show that the circumference of the circle $x^{2}+y^{2}=1$ is $2\pi$.
\end{enumerate}

A nonnegative function is a probability density function if it satisfies the following definition: $\int_{-\infty}^{\infty}f(t)\,dt=1$. The probability that a random variable $x$ lies between $a$ and $b$ is given by $P(a\le x\le b)=\int_{a}^{b}f(t)\,dt$.

\begin{enumerate}[resume]
\item Show that $f(x)=\begin{cases}0 & \text{if }x<0, \\ 7e^{-7x} & \text{if }x\ge 0\end{cases}$ is a probability density function.
\item Find the probability that $x$ is between $0$ and $0.3$. (Use the function defined in the preceding problem.) Use four-place decimal accuracy.
\end{enumerate}

\end{document}
